\chapter{Segurança da Informação}

\section{Princípios e modelo de proteção}

\begin{teoria}[propriedades de segurança]
\begin{itemize}
\item \textbf{Confidencialidade}: informação acessível somente a entidades autorizadas.
\item \textbf{Integridade}: proteção contra alteração ou destruição não autorizada e preservação de exatidão/completude.
\item \textbf{Disponibilidade}: acesso e uso quando necessário.
\item \textbf{Autenticidade}: confiança na identidade/origem.
\item \textbf{Responsabilização (accountability)}: capacidade de associar ações a entidades e manter evidência.
\item \textbf{Não repúdio}: mecanismos que dificultam negar autoria/participação em uma transação, dentro do modelo jurídico e criptográfico aplicável.
\end{itemize}
\end{teoria}

Segurança eficaz trabalha em \termo{defesa em profundidade}: controles preventivos, detectivos, corretivos e compensatórios em pessoas, processos e tecnologia. Nenhum controle isolado deve ser tratado como garantia absoluta.

\section{Política e gestão de segurança}

Política define direção e princípios. Normas/padrões tornam requisitos mais específicos; procedimentos detalham execução; guias orientam boas práticas. Em governança, documentos precisam de dono, escopo, aprovação, comunicação, revisão e mecanismos de conformidade.

Privilégio mínimo concede apenas acessos necessários; need-to-know limita acesso conforme necessidade; segregação de funções reduz possibilidade de uma pessoa iniciar, aprovar e ocultar operação crítica.

\section{Segurança física e lógica}

Controles físicos: perímetro, vigilância, controle de entrada, ambiente, energia, incêndio e descarte seguro. Controles lógicos: autenticação, autorização, criptografia, hardening, firewall, EDR, logging, patching e segmentação. Um data center pode ser logicamente seguro e fisicamente vulnerável — e vice-versa.

\section{Autenticação e autorização}

Fatores clássicos: algo que o usuário sabe, possui ou é. MFA exige fatores de categorias distintas; duas senhas não formam autenticação multifator. Tokens podem ser físicos ou lógicos. Biométricos têm taxas de falso aceite e falsa rejeição e não são ``segredos'' facilmente substituíveis.

Autenticação responde \emph{quem é?}; autorização, \emph{o que pode fazer?}. Modelos incluem RBAC (papéis), ABAC (atributos/políticas) e DAC/MAC em outros contextos.

\section{Malware e golpes}

\begin{table}[ht]
\centering
\small
\begin{tabularx}{\textwidth}{l X}
\toprule
Ameaça & Característica \\
\midrule
Vírus & anexa-se/associa-se a conteúdo hospedeiro e depende de execução para propagação. \\
Worm & propaga-se autonomamente, tipicamente explorando rede/vulnerabilidades. \\
Trojan & aparenta função legítima enquanto executa ação maliciosa; não implica autorreplicação. \\
Rootkit & busca ocultar presença e manter acesso privilegiado. \\
Spyware & coleta informações sem consentimento adequado. \\
Adware & exibe publicidade e pode incorporar rastreamento indesejado. \\
Backdoor & caminho oculto/alternativo de acesso, legítimo ou malicioso conforme contexto. \\
Bot/botnet & host controlado remotamente; conjunto coordenado forma botnet. \\
Ransomware & bloqueia/cripta dados ou extorque por ameaça de vazamento; variantes modernas combinam exfiltração e indisponibilidade. \\
\bottomrule
\end{tabularx}
\end{table}

Phishing usa engenharia social para induzir vítima; spear phishing é direcionado. Pharming redireciona tráfego para destino fraudulento por manipulação de resolução/roteamento/configuração. Hoax é boato fraudulento. Advance fee fraud promete benefício futuro em troca de pagamento antecipado.

\section{Ataques em rede e aplicação}

\subsection{Reconhecimento e interceptação}
Scan identifica hosts, portas e serviços. Sniffing captura tráfego. Spoofing falsifica identidade/origem em algum nível, como IP, e-mail ou ARP. Brute force tenta credenciais/chaves por repetição; credential stuffing reutiliza pares vazados.

\subsection{DoS e DDoS}
Negação de serviço esgota recurso ou explora falha para indisponibilizar serviço. DDoS distribui origem por múltiplos agentes. Mitigação pode envolver rate limiting, filtragem, CDN/anti-DDoS, capacidade, anycast, proteção na operadora e desenho resiliente.

\subsection{SQL injection e buffer overflow}
SQL injection ocorre quando entrada não confiável altera estrutura semântica de consulta. Controles centrais: consultas parametrizadas/prepared statements, validação, privilégio mínimo e tratamento seguro de erros. Buffer overflow ocorre quando escrita excede região prevista, podendo corromper memória e controle de fluxo; mitigadores incluem linguagens seguras, bounds checking, ASLR, DEP/NX e stack canaries.

\section{DMZ, firewall, IDS, IPS e proxy}

DMZ isola serviços expostos da rede interna. Firewall aplica política de filtragem; stateful firewall acompanha estado de conexões. WAF inspeciona tráfego de aplicação web. IDS detecta e alerta; IPS atua inline e pode bloquear. Proxy intermedeia conexões; reverse proxy fica diante de servidores e pode terminar TLS, balancear e aplicar políticas.

\begin{cebraspe}
IDS não é sinônimo de IPS. Firewall não detecta toda ameaça. WAF não substitui correção de vulnerabilidade. DMZ não é ``rede sem proteção'', mas uma zona com controles e exposição deliberadamente limitada.
\end{cebraspe}

\section{VPN}

VPN cria canal lógico protegido sobre rede não confiável. Pode ser site-to-site ou acesso remoto. IPsec atua em camada de rede, com AH/ESP conforme arquitetura; TLS sustenta várias VPNs de aplicação/transporte. VPN não garante que endpoint esteja íntegro: dispositivo comprometido pode entrar no túnel legitimamente.

\section{Criptografia}

\subsection{Simétrica e assimétrica}
Criptografia simétrica usa segredo compartilhado e é eficiente para grandes volumes. Assimétrica usa par de chaves e permite casos como troca de chaves e assinaturas. Em sistemas reais, arquiteturas híbridas usam assimétrica para autenticação/acordo e simétrica para dados.

AES é cifra simétrica moderna. RSA e criptografia de curvas elípticas são assimétricas. Hash criptográfico produz resumo de tamanho fixo e deve resistir, conforme propriedade, a preimagem, segunda preimagem e colisão. Hash não é criptografia reversível.

\subsection{Assinatura digital}
Em modelo simplificado, o signatário calcula hash e gera assinatura com operação vinculada à sua chave privada; verificação usa a chave pública. Assinatura busca autenticidade, integridade e não repúdio técnico/jurídico conforme infraestrutura. \textbf{Não fornece confidencialidade por si.}

\section{PKI e certificados}

PKI associa identidades a chaves públicas por certificados. Autoridade Certificadora emite/assina certificados; Autoridade de Registro valida identidades conforme política. Cadeia de confiança liga certificado final a âncora confiável. Revogação pode ser consultada por CRL/OCSP conforme ecossistema.

\section{Honeypots e honeynets}

Honeypot é recurso deliberadamente atrativo/monitorado para detectar, estudar ou desviar atividade maliciosa. Honeynet é rede de honeypots/ambiente de engodo. Exigem isolamento e monitoramento: um honeypot comprometido não pode virar plataforma de ataque a terceiros.

\section{Gestão de riscos}

Risco relaciona ativos/objetivos, ameaças, vulnerabilidades, probabilidade e impacto. Tratamentos típicos: evitar, mitigar/modificar, transferir/compartilhar ou aceitar. Risco inerente é o risco antes de controles; residual permanece após controles.

\begin{exemplo}[quantificação simplificada]
Se um incidente tem perda esperada de R\$ 500 mil por ocorrência e frequência média estimada de 0,2/ano, a perda anual esperada simplificada é $500000\times0{,}2=R\$100000$. Esse cálculo apoia decisão, mas não captura caudas, incerteza e efeitos não financeiros.
\end{exemplo}

\section{ISO/IEC 27001, 27002 e 27005}

ISO/IEC 27001 estabelece requisitos para sistema de gestão de segurança da informação (SGSI). ISO/IEC 27002 fornece orientações de controles. ISO/IEC 27005 trata gestão de riscos de segurança da informação. A prova pode exigir a diferença entre \textbf{requisito de sistema de gestão}, \textbf{catálogo/orientação de controles} e \textbf{processo de risco}.

\section{Continuidade de negócio}

Análise de impacto no negócio (BIA) identifica processos críticos, dependências e impactos ao longo do tempo. RTO é tempo-alvo para restabelecimento; RPO é ponto no tempo até o qual dados devem ser recuperados, equivalendo à perda temporal tolerável. MTPD/MAO refere-se ao período máximo tolerável de interrupção conforme terminologia adotada.

Planos precisam ser testados. Backup sem restauração testada é evidência insuficiente de recuperabilidade.

O edital cita ISO/IEC 15999 e também, em governança, ISO 22301. Como 15999 é referência histórica de continuidade, estude a evolução conceitual e priorize os princípios de SG de continuidade: política, BIA, estratégia, planos, exercícios, avaliação e melhoria.

\section{Mapa mental}

\mapamental{Segurança}{
 child {node[concept]{Princípios} child {node[concept]{CIA}} child {node[concept]{accountability}}}
 child {node[concept]{IAM} child {node[concept]{MFA}} child {node[concept]{RBAC}}}
 child {node[concept]{Ameaças} child {node[concept]{malware}} child {node[concept]{phishing}}}
 child {node[concept]{Rede} child {node[concept]{DMZ}} child {node[concept]{IDS/IPS}}}
 child {node[concept]{Cripto} child {node[concept]{AES/RSA}} child {node[concept]{PKI}}}
 child {node[concept]{Gestão} child {node[concept]{27001}} child {node[concept]{risco/BCP}}}
}

\section{Questões por nível}

\begin{questao}{\nivel{N1} 13.1}
Qual controle tem como função primária detectar tráfego suspeito e gerar alerta, sem necessariamente estar em posição de bloqueio inline?
\begin{enumerate}[label=\Alph*)]
\item IDS
\item IPS obrigatoriamente bloqueante
\item RAID
\item NAT
\item PKI
\end{enumerate}
\end{questao}
\begin{solucao}{13.1}
\textbf{Gabarito: A.} IDS é tipicamente detectivo. IPS opera inline para prevenir/bloquear conforme política. RAID é armazenamento; NAT traduz endereços; PKI gerencia confiança/chaves/certificados.
\end{solucao}

\begin{questao}{\nivel{N2} 13.2}
Um usuário autentica-se com senha e PIN enviados pelo mesmo aplicativo. Isso é, por definição, MFA?
\begin{enumerate}[label=\Alph*)]
\item Sim, pois há duas sequências secretas.
\item Sim, desde que cada uma tenha oito caracteres.
\item Não necessariamente; senha e PIN pertencem ao mesmo fator de conhecimento.
\item Não, porque MFA só admite biometria e certificado.
\item Sim, se o canal usar TLS.
\end{enumerate}
\end{questao}
\begin{solucao}{13.2}
\textbf{Gabarito: C.} MFA combina fatores de categorias diferentes. A/B confundem número de credenciais com fatores; D restringe indevidamente; E protege canal, mas não transforma dois conhecimentos em fatores distintos.
\end{solucao}

\begin{questao}{\nivel{N3} 13.3}
Qual propriedade uma assinatura digital oferece diretamente, mas não a criptografia simétrica de dados por si só?
\begin{enumerate}[label=\Alph*)]
\item Confidencialidade obrigatória.
\item Prova vinculada à chave privada do signatário, apoiando autenticidade e integridade.
\item Compressão sem perdas.
\item Disponibilidade do serviço.
\item Anonimização de dados pessoais.
\end{enumerate}
\end{questao}
\begin{solucao}{13.3}
\textbf{Gabarito: B.} Assinatura vincula assinatura à chave privada e permite verificação com chave pública. Ela não cifra necessariamente o conteúdo (A); C, D e E são propriedades distintas.
\end{solucao}

\begin{questao}{\nivel{N4} 13.4}
Auditoria encontra administradores capazes de alterar registros de negócio e apagar os próprios logs de administração. Qual é o principal problema de controle?
\begin{enumerate}[label=\Alph*)]
\item ausência de RAID 6.
\item quebra de segregação e fragilidade de accountability/integridade da trilha de auditoria.
\item uso de criptografia assimétrica.
\item excesso de disponibilidade.
\item inexistência de NAT.
\end{enumerate}
\end{questao}
\begin{solucao}{13.4}
\textbf{Gabarito: B.} Quem realiza ação crítica não deve poder destruir evidência sem controle independente. A/C/E são tecnologias não relacionadas ao problema central; D é sem sentido no cenário.
\end{solucao}

\begin{discursiva}[ransomware]
Um órgão sofreu ransomware com exfiltração, criptografia de servidores e comprometimento de credenciais administrativas. Backups online estavam acessíveis pelo mesmo domínio. Em até 30 linhas, apresente falhas prováveis, medidas imediatas de resposta e controles estruturais para reduzir recorrência e impacto.
\end{discursiva}
